\chapter{Wstęp}
\section{Wprowadzenie}
Sieci 5G Core są rozwijane w modelu \emph{cloud native}, w którym funkcje sieciowe (CNF) są uruchamiane jako kontenery i zarządzane przez platformy orkiestracyjne, takie jak Kubernetes~\cite{3gpp23_501,kubernetesDocs}. Taka architektura zwiększa elastyczność i skalowalność, ale jednocześnie podnosi ryzyko operacyjne związane z częstymi wdrożeniami nowych wersji oprogramowania.

W środowiskach telekomunikacyjnych nawet krótkotrwała regresja może prowadzić do zakłóceń sygnalizacji, wzrostu opóźnień oraz degradacji jakości usług dla dużej liczby abonentów~\cite{3gpp23_502}. W związku z tym konieczne jest stosowanie metod wdrożeniowych, które pozwalają testować nowe wersje funkcji sieciowych na realnym ruchu, ale bez wpływu na użytkowników końcowych.

\section{Opis problemu}
W klasycznym modelu wdrożeń aktualizacja funkcji 5G Core bywa wykonywana bezpośrednio na ścieżce produkcyjnej. Przy dużej skali obciążenia i współzależnościach między funkcjami SBA (\emph{Service-Based Architecture}) taka praktyka zwiększa prawdopodobieństwo:
\begin{itemize}
    \item niekontrolowanego wzrostu błędów sygnalizacyjnych,
    \item przeciążenia wybranych komponentów,
    \item uszkodzenia lub nadpisania stanu w repozytoriach danych abonentów,
    \item trudności w szybkim wycofaniu zmiany.
\end{itemize}

Szczególnie wrażliwym obszarem jest zależność funkcji kontrolnych (np. SMF) od danych przechowywanych poza samą funkcją, m.in. w UDR/UDSF. Jeżeli nowa wersja wykonuje operacje zapisu na produkcyjnym stanie, test wdrożeniowy przestaje być bezpieczny.

\section{Cel i zakres pracy}
Głównym celem pracy jest opracowanie kompletnej architektury \emph{Shadow Deploy} dla funkcji SMF w 5G Core, która:
\begin{itemize}
    \item wykorzystuje Istio/Envoy do bezinwazyjnego klonowania ruchu,
    \item izoluje warstwę danych wersji testowej od warstwy produkcyjnej,
    \item wspiera automatyczną ocenę jakości na podstawie metryk i progów SLO.
\end{itemize}

Zakres pracy obejmuje przygotowanie artefaktów infrastrukturalnych i konfiguracji:
\begin{itemize}
    \item definicji środowiska Kubernetes i namespace z aktywną iniekcją sidecarów,
    \item manifestów SMF v1 (stabilna) i SMF v2 (shadow),
    \item reguł routingu Istio (\texttt{DestinationRule}, \texttt{VirtualService}),
    \item konfiguracji walidacyjnej opartej na Prometheus i Iter8,
    \item opcjonalnego komponentu porównującego odpowiedzi (Diffy).
\end{itemize}

\section{Teza i pytania badawcze}
W pracy przyjęto tezę, że zastosowanie architektury Shadow Deploy w 5G Core umożliwia istotne ograniczenie ryzyka wdrożeniowego przy zachowaniu realnej reprezentatywności testów.

Z tezy wynikają trzy pytania badawcze:
\begin{enumerate}
    \item Czy mirrorowanie ruchu na poziomie Service Mesh pozwala wiarygodnie oceniać wersję kandydacką SMF bez wpływu na użytkownika?
    \item Czy separacja stanu danych (UDR produkcyjny vs UDR shadow) eliminuje ryzyko niepożądanego zapisu podczas testów?
    \item Czy metryki z Envoy i Prometheus są wystarczające do automatycznej decyzji \emph{pass/fail} dla nowej wersji?
\end{enumerate}

\section{Struktura pracy}
Praca została podzielona na sześć rozdziałów:
\begin{itemize}
    \item Rozdział 1 definiuje kontekst problemu, cele i zakres.
    \item Rozdział 2 opisuje wymagania, architekturę logiczną i artefakty IaC.
    \item Rozdział 3 przedstawia implementację środowiska Shadow Deploy.
    \item Rozdział 4 opisuje metodykę walidacji i scenariusze testowe.
    \item Rozdział 5 prezentuje wyniki eksperymentów i ich analizę.
    \item Rozdział 6 zawiera podsumowanie oraz kierunki dalszego rozwoju.
\end{itemize}